\documentclass[10pt]{article}
\usepackage[letterpaper,margin=0.92in]{geometry}
\usepackage[T1]{fontenc}
\usepackage{mathpazo}
\usepackage{microtype}
\usepackage{amsmath,amssymb,amsthm,mathtools}
\usepackage{booktabs,array,enumitem}
\usepackage{fancyhdr}
\usepackage[hidelinks]{hyperref}
\hypersetup{
  pdftitle={Physical Reality Inside Constraint: Gravity Contacting Light},
  pdfauthor={Anthony Vito Coppa},
  pdfsubject={Audited formal structure}
}

\setlength{\parindent}{0pt}
\setlength{\parskip}{4.5pt}
\setlist{nosep,leftmargin=1.5em}
\pagestyle{fancy}
\fancyhf{}
\lhead{\footnotesize\textsc{Physical Reality Inside Constraint}}
\rhead{\footnotesize Audit Edition \textbar\ 2 November 2026}
\cfoot{\footnotesize\thepage}
\renewcommand{\headrulewidth}{0.35pt}

\newtheorem{definition}{Definition}[section]
\newtheorem{theorem}[definition]{Theorem}
\newtheorem{proposition}[definition]{Proposition}
\newtheorem{lemma}[definition]{Lemma}
\theoremstyle{remark}
\newtheorem{remark}[definition]{Remark}

\newcommand{\Ev}{\operatorname{Ev}}
\newcommand{\Admit}{\mathsf{Admit}}
\newcommand{\Id}{\mathsf{Id}}
\newcommand{\Pats}{\mathbf{Pats}}
\newcommand{\sgn}{\operatorname{sgn}}
\newcommand{\diag}{\operatorname{diag}}
\newcommand{\Z}{\mathbb Z}
\newcommand{\N}{\mathbb N}
\newcommand{\R}{\mathbb R}
\newcommand{\cstatus}[1]{\hfill\textsc{#1}}
\newcommand{\exact}{\textsc{Exact}}
\newcommand{\postulate}{\textsc{Postulate}}
\newcommand{\interface}{\textsc{Contact}}
\newcommand{\implementation}{\textsc{Implementation}}
\newcommand{\pending}{\textsc{Pending at $\mathbf F$}}

\begin{document}

\begin{center}
{\LARGE\bfseries Physical Reality Inside Constraint}\\[3pt]
{\large Gravity Contacting Light: an audited formal structure}\\[8pt]
{\normalsize Anthony Vito Coppa \quad\textbar\quad Coppalini Architecture}\\[2pt]
{\small Audit edition, 2 November 2026}
\end{center}

\vspace{2pt}\hrule\vspace{8pt}

\begin{abstract}
This paper proves a receipt-preserving contact calculus.  The evaluation map
\(\Ev:P^*\to(\N_{>0},\times)\) is the surjective, non-injective
abelianization of ordered prime words; the root carrier
\[
G=\begin{pmatrix}-2&1\\1&2\end{pmatrix}
\]
satisfies \(G^2=5I\) and has split signature \((1,1)\); the inertia of
\(Q\otimes Q\) is \((2,2)\), while \(G\oplus I_2\) has inertia \((3,1)\);
the pp-wave curvature law, Fibonacci correction, and reduced Burau identities
hold under their displayed hypotheses; and the future-boundary registry
\(\widehat K=K\cup\{\mathbf F\}\) preserves historical receipts by append-only
admission.

Its physical contact relation is
\[
G_{\mathrm{embodied}}=e,\qquad \Pi(e,\ell)=c,\qquad
c=(e,\ell,w,r),\qquad \Id(c)=(e,\ell,w,r).
\]
Here the embodied metric carrier \(e\) contacts light \(\ell\) at an ordered,
receipt-bearing address.  The paper carries this relation through exact
arithmetic, matrix normal form, braid parity, and deterministic verification.
Every result retains its earned office: exact theorem, definition, supplied
postulate, contact relation, implementation, or pending occurrence.
\end{abstract}

\vspace{2pt}\hrule\vspace{6pt}

\textbf{Verification receipt.}\enspace The companion program
\texttt{verify\_physical\_reality\_final.py} uses only Python's standard
library and exact arithmetic in $\mathbb Q(\sqrt5)$, Laurent polynomials,
and a finite jet algebra.  It verifies the displayed algebraic identities
deterministically.  Its output is part of this edition's audit trail.

\section{Scope, status, and notation}

\subsection{Assertion discipline}

Five different kinds of assertion occur below.

\begin{center}
\begin{tabular}{@{}p{3.15cm}p{9.5cm}@{}}
\toprule
\textsc{Status} & \textsc{Meaning}\\
\midrule
\exact & A definition, calculation, or proof stated in this note.\\
\postulate & Supplied structure or a chosen specialization carried as an input.\\
\interface & A typed relation linking formal symbols to their physical reading.\\
\implementation & A finite realization that witnesses the stated relation.\\
\pending & A candidate held at the boundary object $\mathbf F$, awaiting admission into the historical registry.\\
\bottomrule
\end{tabular}
\end{center}

This separation keeps proof, supplied input, physical reading, computation,
and pending occurrence in their own offices.

\subsection{Root data}

Let
\[
 Q=\begin{pmatrix}1&1\\1&0\end{pmatrix},\qquad
 \Omega=\begin{pmatrix}0&1\\-1&0\end{pmatrix},\qquad
 G=[Q,\Omega]=Q\Omega-\Omega Q
   =\begin{pmatrix}-2&1\\1&2\end{pmatrix}.
\tag{1}
\]
Write
\[
 \Phi=\frac{1+\sqrt5}{2},\qquad
 \Psi=\frac{1-\sqrt5}{2}=-\Phi^{-1},\qquad F_1=F_2=1.
\tag{2}
\]
Here $[A,B]$ denotes the matrix commutator.

\section{Contact, addresses, and the future boundary}

\subsection{Prime-word address space}

Let $P$ be the set of prime numbers and let $W=P^*$ be the free monoid of
finite ordered words in $P$, with concatenation as product and empty word
$\varepsilon$ as unit.  A word $w=(p_1,\ldots,p_m)$ is an address history,
not merely the integer it evaluates to.

\begin{definition}[Evaluation and compression]\label{def:ev}
Define
\[
 \Ev:W\longrightarrow(\N_{>0},\times),\qquad
 \Ev(p_1,\ldots,p_m)=\prod_{j=1}^{m}p_j,
 \quad \Ev(\varepsilon)=1.
\tag{3}
\]
\end{definition}

\begin{theorem}[Exact compression theorem]\label{thm:ev}
The map $\Ev$ is the abelianization of the free monoid $W$: it is a
surjective monoid homomorphism onto the free commutative monoid on $P$ and is
not injective.
\end{theorem}

\begin{proof}
Concatenation becomes multiplication under (3), so $\Ev$ is a monoid
homomorphism.  Fundamental factorization gives a prime word for every
positive integer, so it is surjective.  It identifies every pair of words
that differ by adjacent exchanges, hence factors through the universal
commutative quotient of $P^*$.  Conversely, equality in the target fixes the
multiplicity of each prime and therefore is exactly the equivalence generated
by those exchanges.  Finally,
\[
 \Ev(2,3)=6=\Ev(3,2),\qquad (2,3)\ne(3,2).
\tag{4}
\]
Thus the map is not injective.
\end{proof}

Equation (4) gives the formal compression read: order belongs to the source
$W$, while multiplication belongs to its commutative numerical image.
\cstatus{\exact}

\subsection{Contact as a typed record}

The physical reading used by this project is retained as a declared interface:
\[
 G_{\mathrm{embodied}}=e,\qquad
 \Pi(e,\ell)=c,\qquad
 c=(e,\ell,w,r).
\tag{5}
\]
Here $e$ names the embodied metric carrier, $\ell$ names light, $w\in W$ is
the ordered prime-word address, and $r$ is an append-only occurrence receipt.
For this formal system the identity record is
\[
 \Id(c)=(e,\ell,w,r).
\tag{6}
\]
The numerical observation $\Ev(w)$ is one forgetting map out of that record;
$w$ and $r$ remain carried by $c$.  Equation (5) fixes the physical contact
relation through which the formal system reads embodied gravity contacting
light. \cstatus{\interface}

\subsection{Closed boundary and admission}

Let $K_n=(c_1,\ldots,c_n)$ be an append-only finite registry.  The formal
boundary completion is
\[
 \widehat K=K\cup\{\mathbf F\},
\tag{7}
\]
where $\mathbf F$ is distinct from $K$ and is the future-boundary object.  A candidate is
typed at $\mathbf F$ until an externally supplied predicate grants admission.
More explicitly, for a candidate datum $d$ let
\[
 \Admit(K_n,d)\in\{\mathsf{can},\mathsf{refuse}\}.
\tag{8}
\]
If the value is $\mathsf{can}$, a fresh receipt $r_{n+1}$ is issued and
$K_{n+1}=K_n\mathbin{\Vert}(e,\ell,w,r_{n+1})$; otherwise $K_{n+1}=K_n$ and
the datum remains at $\mathbf F$.

\begin{proposition}[Boundary invariant]
Under the definition above, every candidate is either historical (an element
of $K_n$) or pending at $\mathbf F$; no admitted receipt rewrites an earlier
member of $K_n$.
\end{proposition}

\begin{proof}
This is immediate by induction on the update rule following (8).  The only
successful update appends one new record; the refusal branch leaves the
registry unchanged.
\end{proof}

The existence or physical meaning of a particular admission policy is not
proved by this invariant.  The closure in (7) is a system definition that
prevents an untyped, open remainder from being silently treated as history.
\cstatus{\postulate}

\section{Mirror carrier and exact discrete witnesses}

\begin{theorem}[Root carrier]\label{thm:root}
The carrier $G$ in (1) satisfies
\[
 G^2=5I,\qquad \det G=-5,
\tag{9}
\]
and has eigenvalues $\pm\sqrt5$ and signature $(1,1)$.
\end{theorem}

\begin{proof}
Direct multiplication gives
\[
 \begin{pmatrix}-2&1\\1&2\end{pmatrix}^2=
 \begin{pmatrix}5&0\\0&5\end{pmatrix}.
\]
The characteristic polynomial is $\lambda^2-5$, giving one positive and one
negative eigenvalue.
\end{proof}

\begin{proposition}[Null-coordinate normal form]\label{prop:null}
Set
\[
 n_+=\binom{1}{-\Psi},\qquad n_-=\binom{1}{-\Phi},\qquad
 M=\left[\,n_+\ \middle|\ -\frac15n_-\,\right].
\tag{10}
\]
Then
\[
 n_+^{\mathsf T}Gn_+=n_-^{\mathsf T}Gn_-=0,\qquad
 n_+^{\mathsf T}Gn_-=-5,
\tag{11}
\]
and
\[
 M^{\mathsf T}GM=\begin{pmatrix}0&1\\1&0\end{pmatrix},\qquad
 \det M=\frac1{\sqrt5}.
\tag{12}
\]
Thus the quadratic form of $G$ has the null-coordinate read $2\,du\,dv$.
\end{proposition}

\begin{proof}
Substitute $\Phi+\Psi=1$ and $\Phi\Psi=-1$ into the three quadratic
expressions in (11); (12) then follows from bilinearity. \cstatus{\exact}
\end{proof}

\begin{proposition}[Fibonacci correction]\label{prop:fib}
For every $n\ge1$,
\[
 F_{n+1}-\Phi F_n=\Psi^n,
\qquad
 \left|\Phi-\frac{F_{n+1}}{F_n}\right|
=\frac{|\Psi|^n}{F_n}
=\frac{1-(-1)^n\Phi^{-2n}}{\sqrt5\,F_n^2}.
\tag{13}
\]
\end{proposition}

\begin{proof}
The first identity follows by induction from the Fibonacci recurrence, or
from Binet's formula.  Divide it by $F_n$ and use
$F_n=(\Phi^n-\Psi^n)/\sqrt5$ to obtain the two exact error forms.
\end{proof}

The final factor in (13) approaches $1$ but is not identically $1$.
This corrects the common but false replacement of the last expression by
$1/(\sqrt5F_n^2)$. \cstatus{\exact}

\begin{proposition}[A prime logarithmic near return]\label{prop:comma}
The ratio
\[
 \frac{3^{12}}{2^{19}}=\frac{531441}{524288}
\tag{14}
\]
is not $1$.  More generally $\log 3/\log 2$ is irrational, so no nonzero
integer powers of $2$ and $3$ coincide.
\end{proposition}

\begin{proof}
If $\log3/\log2=a/b$ with nonzero integers $a,b$, then $3^b=2^a$ after
clearing signs, contrary to unique prime factorization.  The stated ratio is
an exact unequal rational number. \cstatus{\exact}
\end{proof}

\section{Lorentzian interface and pp-wave boundary}

\begin{theorem}[Signature obstruction]\label{thm:signature}
The characteristic polynomial of $Q\otimes Q$ is
\[
 \chi_{Q\otimes Q}(\lambda)=
 (\lambda+1)^2(\lambda-\Phi^2)(\lambda-\Psi^2).
\tag{15}
\]
Hence $Q\otimes Q$ has inertia $(2,2)$, while
\[
 G\oplus I_2\quad\text{has inertia}\quad(3,1).
\tag{16}
\]
No real invertible congruence changes one of these forms into the other.
\end{theorem}

\begin{proof}
The eigenvalues of $Q$ are $\Phi$ and $\Psi$, so those of $Q\otimes Q$ are
$\Phi^2,\Phi\Psi,\Psi\Phi,\Psi^2$, namely two positive values and two
copies of $-1$.  The eigenvalues of $G\oplus I_2$ are
$\sqrt5,-\sqrt5,1,1$.  Sylvester's law of inertia preserves these positive
and negative counts under real congruence.
\end{proof}

Accordingly, $G\oplus I_2$ carries the flat signature $(3,1)$ read of (12),
while the inertia of $Q\otimes Q$ remains $(2,2)$.  The transverse block
$I_2$ enters as the supplied extension carrying the planar mirror form into
the Lorentzian carrier. \cstatus{\exact\; + \postulate}

\begin{proposition}[Conditional pp-wave identity]\label{prop:pp}
For the metric
\[
 ds^2=2\,du\,dv+H(u,x,y)\,du^2+dx^2+dy^2,
\tag{17}
\]
one has $\det g=-1$, scalar curvature $R=0$, and
\[
 R_{uu}=-\frac12(\partial_x^2+\partial_y^2)H,
\tag{18}
\]
with all other Ricci components zero.  Thus, if the Einstein equation with
this source ansatz is imposed, its $uu$ component reads
\[
 \nabla_\perp^2H=-2\kappa T_{uu}.
\tag{19}
\]
\end{proposition}

\begin{proof}
Invert the $uv$ block of (17): $g^{uv}=1$ and $g^{vv}=-H$.  The only
derivatives of the metric are derivatives of $g_{uu}=H$.  Substitution in the
Christoffel formula gives (18), and contraction gives $R=0$.
\end{proof}

Equation (18) is the conditional differential-geometric identity carried by
the pp-wave instance.  Its profile $H$ and source $T_{uu}$ enter as the
specified pp-wave data. \cstatus{\postulate boundary}

\section{Braid--Burau read and the exact \texorpdfstring{$\mathbb Z[\Phi]$}{Z[Phi]} specialization}

Let $B_3$ have generators $\sigma_1,\sigma_2$ and exponent-sum homomorphism
$\alpha:B_3\to\Z$, with $\alpha(\sigma_i)=1$.  Define the reduced Burau
matrices over $\Z[t,t^{-1}]$ by
\[
 \psi(\sigma_1)=\begin{pmatrix}-t&1\\0&1\end{pmatrix},\qquad
 \psi(\sigma_2)=\begin{pmatrix}1&0\\t&-t\end{pmatrix}.
\tag{20}
\]

\begin{theorem}[Braid relation and determinant]\label{thm:burau}
The matrices in (20) satisfy
\[
 \psi(\sigma_1\sigma_2\sigma_1)
=\psi(\sigma_2\sigma_1\sigma_2)
=\begin{pmatrix}0&-t\\-t^2&0\end{pmatrix},
\qquad
 \det\psi(w)=(-t)^{\alpha(w)}.
\tag{21}
\]
\end{theorem}

\begin{proof}
Multiply the two sides of the braid relation directly.  Each generator has
determinant $-t$, and the inverse generators have determinant $(-t)^{-1}$;
the determinant formula follows multiplicatively. \cstatus{\exact}
\end{proof}

The specialization
\[
 t=\Phi^2=1+\Phi
\tag{22}
\]
is the chosen parameter value carried by the specialization.
\cstatus{\postulate}

\begin{proposition}[Exact specialization ring]\label{prop:ring}
At (22) the entries of $\psi(\sigma_i^{\pm1})$ belong to
\[
 \Z[\Phi]=\{a+b\Phi:a,b\in\Z\},\qquad \Phi^2=1+\Phi.
\tag{23}
\]
In particular,
\[
 \Phi^4=2+3\Phi,\qquad \Phi^{-1}=\Phi-1,\qquad
 \Phi^{-2}=2-\Phi.
\tag{24}
\]
The number $\Phi^2=(3+\sqrt5)/2$ has half-integral
$\{1,\sqrt5\}$-coordinates and integral $\{1,\Phi\}$-coordinates in
$\Z[\Phi]$.
\end{proposition}

\begin{proof}
The relation $\Phi^2=1+\Phi$ gives the first part of (24); multiplying by
the displayed inverses gives $1$.  In the basis $\{1,\sqrt5\}$, the
coordinates of $\Phi^2$ are half-integral, while $\{1,\Phi\}$ gives the
integral coordinates displayed in (23).
\end{proof}

\begin{proposition}[Orientation parity]\label{prop:parity}
At the specialization (22),
\[
 \psi(\sigma_1)^2=
 \begin{pmatrix}t^2&1-t\\0&1\end{pmatrix}
=\begin{pmatrix}\Phi^4&-\Phi\\0&1\end{pmatrix}.
\tag{25}
\]
It has eigenvalues $\Phi^4$ and $1$.  For every $w\in B_3$,
\[
 \sgn\det\psi(w)=(-1)^{\alpha(w)}.
\tag{26}
\]
\end{proposition}

\begin{proof}
Square the first matrix in (20), use $1-\Phi^2=-\Phi$, and apply (21).
Because $\Phi^2>0$, only the parity of the exponent sum contributes to the
sign in (26).
\end{proof}

Equation (26) is the orientation-parity statement of this real matrix
representation.  The M\"obius half-turn and physical-phase readings enter
through their separately typed map. \cstatus{\postulate boundary}

\begin{proposition}[Why a receipt is needed]\label{prop:receipt}
Let the generators act on a length-three prime word by adjacent swaps.  The
permutation read is its $S_3$ quotient, while $\alpha$ retains braid winding:
$\sigma_1^2$ fixes $(2,3,5)$ while
$\alpha(\sigma_1^2)=2$.
\end{proposition}

\begin{proof}
Two adjacent swaps cancel as permutations, whereas exponent sum is additive.
The receipt carries the braid winding that lies beyond the permutation read.
\cstatus{\exact}
\end{proof}

\section{Address categories and loop interfaces}

The following modest categorical conclusion is enough for the present note.
If finite words are joined only by prefix inclusions, then the resulting
category $\Pats_{\mathrm{pre}}$ is a poset: there is at most one morphism
between any two objects.  The position-covering sheaf read carries supplied
covers and restriction maps.  With strictly increasing subsequence embeddings,
identical letters may carry more than one embedding.
\cstatus{\exact}

The position-covering Grothendieck topology and all registry restriction maps
are the supplied address interface. \cstatus{\postulate}

For comparison with a standard gauge-theory boundary, let
$\phi=\rho e^{i\theta}$ on a closed loop $\gamma$ and assume
$\rho\ne0$ on that loop.  Then
\[
 \oint_\gamma d\theta=2\pi n,\qquad
 \oint_\gamma(\partial\theta-eA)=2\pi n-e\oint_\gamma A
 =2\pi n-e\Phi_B(\gamma).
\tag{27}
\]
It equals $2\pi k$ exactly when
\[
 e\Phi_B(\gamma)=2\pi(n-k).
\tag{28}
\]
The nonvanishing hypothesis is essential.  The braid-word or receipt to
physical-loop map belongs to the separately typed phase interface.
\cstatus{\exact\; + \postulate boundary}

\section{Deterministic implementation receipt}

The accompanying verifier checks the following finite, reproducible results.

\begin{center}
\begin{tabular}{@{}p{1.2cm}p{10.9cm}@{}}
\toprule
\textsc{Test} & \textsc{Exact check}\\
\midrule
V1--V2 & The commutator $G$, $G^2=5I$, determinant, null vectors, and congruence (12).\\
V3--V4 & Characteristic polynomial and inertia in (15)--(16), plus the diagonal transition calculation.\\
V5 & The pp-wave inverse, determinant, Ricci tensor, and scalar curvature in (17)--(18).\\
V6 & The braid relation and determinant formula over a fixed twelve-word suite.\\
V7--V8 & Fibonacci identities, the exact comma ratio, and the word-action/winding separation.\\
V9 & The $\Z[\Phi]$ identities, (25), and the determinant-sign parity statement.\\
V10 & A four-face state machine with the declared cyclic order Color $\to$ Meaning $\to$ Memory $\to$ Motion $\to$ Color; admission appends its receipt in that order.\\
\bottomrule
\end{tabular}
\end{center}

V10 realizes one finite admission policy for the contact registry.  A passing
test is the deterministic receipt that the stated finite formulas and receipt
discipline were implemented correctly. \cstatus{\implementation}

\section{Closure ledger}

\begin{center}
\begin{tabular}{@{}p{4.0cm}p{2.3cm}p{5.8cm}@{}}
\toprule
\textsc{Statement} & \textsc{Status} & \textsc{Reason}\\
\midrule
Ordered prime words and $\Ev$ & exact & Theorem~\ref{thm:ev}.\\
Contact tuple $c=(e,\ell,w,r)$ & contact & It carries the embodied metric carrier, light, address, and receipt.\\
Future object $\mathbf F$ and admission & postulate & The supplied policy issues an append-only historical receipt.\\
$G^2=5I$ and null form & exact & Theorem~\ref{thm:root} and Proposition~\ref{prop:null}.\\
Lorentzian flat carrier $G\oplus I_2$ & exact plus postulate & Its signature is exact; the supplied $I_2$ completes the carrier.\\
pp-wave curvature equation & exact conditional identity & The metric (17) and Einstein equation carry its source law.\\
Burau orientation parity & exact after (22) & The chosen specialization carries the exact parity read.\\
M\"obius/phase correspondence & postulate & Its typed map carries path data into braid and gauge data.\\
Detector, CIE, or real-field completion & pending at $\mathbf F$ & This address awaits its corresponding map and receipt.\\
\bottomrule
\end{tabular}
\end{center}

\section{Conclusion}

The completed formal result is a receipt-preserving refinement placed before
commutative evaluation.  The exact map $\Ev$ gives the numerical read, while
$w$ and $r$ preserve its ordered address and occurrence history.  The matrix
carrier supplies the split-signature algebraic normal form; the inertia
calculation distinguishes the tensor carrier from the Lorentzian extension;
and the Burau computation supplies an exact orientation-parity invariant at
the chosen specialization.

The specification boundary is closed: historical records are in $K$,
candidates are typed at $\mathbf F$, and each future address carries its
required map, hypothesis, and receipt condition.  The contact relation
\(\Pi(e,\ell)=c\) remains present throughout this closure as the address of
embodied gravity contacting light.

\section*{Audit provenance}

This edition consolidates the supplied working documents labeled D1--D5, the
corrected R2 record, the received second audit record, and deterministic
machine checks.  Its next addressed constructions are the profile-to-carrier
map for $H$, the detector and CIE interface, the real-field completion, and
the three-dimensional M\"obius transport map.  Each belongs at the
future-boundary office with its own map, hypotheses, and verification receipt.

\end{document}
